% ==============================================================================
% Machine Learning & Statistical Theory Notes — LaTeX Template
% Author: Noted by Md. Inzamamul Lohani (Software Engineering, IIT, DU)
% Based on: Lasso & Ridge Regression Revision Notes Architecture
% Compile: pdflatex template.tex && pdflatex template.tex
% ==============================================================================

\documentclass[11pt,a4paper,oneside]{book}

% ==================== Core Packages ====================
\usepackage[T1]{fontenc}
\usepackage{lmodern}
\usepackage[a4paper,margin=25mm,headheight=15pt]{geometry}
\usepackage{microtype}
\usepackage{amsmath,amssymb,mathtools,bm}
\usepackage{graphicx}
\usepackage{booktabs,tabularx,array}
\usepackage{enumitem}
\usepackage{float}
\usepackage[section]{placeins}
\usepackage{xcolor}
\usepackage[most]{tcolorbox}
\usepackage{listings}
\usepackage{fancyhdr}
\usepackage{titlesec}
\usepackage{hyperref}
\usepackage[nameinlink,noabbrev]{cleveref}
\usepackage{caption}

% ==================== Color Palette ====================
\definecolor{Navy}{HTML}{17324D}        % Primary headings, main emphasis
\definecolor{Blue}{HTML}{2878B5}        % Secondary headings, key idea frame
\definecolor{Teal}{HTML}{168477}        % Sub-headings, intuition frame
\definecolor{Gold}{HTML}{D18B26}        % Warnings, alerts frame
\definecolor{PaleBlue}{HTML}{EDF5FB}    % Key idea background
\definecolor{PaleTeal}{HTML}{EAF6F3}    % Intuition background
\definecolor{PaleGold}{HTML}{FFF6E6}    % Warning background
\definecolor{Ink}{HTML}{263238}         % Body text
\definecolor{Muted}{HTML}{61717D}       % Headers, footers, meta text
\definecolor{CodeBg}{HTML}{F7F8FA}      % Code block background
\definecolor{CodeGreen}{HTML}{236E41}   % Code comments
\definecolor{CodePurple}{HTML}{73378C}  % Code strings

% ==================== Hyperlink Setup ====================
\hypersetup{
  colorlinks=true,
  linkcolor=Navy,
  urlcolor=Blue,
  citecolor=Teal,
  pdftitle={Theory Notes Template},
  pdfauthor={Md. Inzamamul Lohani}
}

% Graphics directory
\graphicspath{{figures/}}

% Paragraph & list spacing
\setlength{\parskip}{5pt}
\setlength{\parindent}{0pt}
\setlist[itemize]{leftmargin=1.5em,itemsep=2pt,topsep=3pt}
\setlist[enumerate]{leftmargin=1.7em,itemsep=3pt,topsep=3pt}
\captionsetup{font=small,labelfont=bf}

% ==================== Styled Callout Boxes ====================
% 1. Key Idea Box (Blue)
\newtcolorbox{keyidea}[1][]{
  enhanced,breakable,
  colback=PaleBlue,colframe=Blue,
  boxrule=.7pt,arc=2mm,
  left=3.5mm,right=3.5mm,top=2.5mm,bottom=2.5mm,
  title=Key idea,fonttitle=\bfseries\color{Navy},
  #1
}

% 2. Intuition Box (Teal)
\newtcolorbox{intuition}[1][]{
  enhanced,breakable,
  colback=PaleTeal,colframe=Teal,
  boxrule=.7pt,arc=2mm,
  left=3.5mm,right=3.5mm,top=2.5mm,bottom=2.5mm,
  title=Intuition,fonttitle=\bfseries\color{Navy},
  #1
}

% 3. Warning / Important Note Box (Gold)
\newtcolorbox{warningbox}[1][]{
  enhanced,breakable,
  colback=PaleGold,colframe=Gold,
  boxrule=.7pt,arc=2mm,
  left=3.5mm,right=3.5mm,top=2.5mm,bottom=2.5mm,
  title=Important note,fonttitle=\bfseries\color{Navy},
  #1
}

% 4. Worked Example Box (Clean border)
\newtcolorbox{worked}[1][]{
  enhanced,breakable,
  colback=white,colframe=Muted!65,
  boxrule=.6pt,arc=2mm,
  left=3.5mm,right=3.5mm,top=2.5mm,bottom=2.5mm,
  title=Worked example,fonttitle=\bfseries\color{Navy},
  #1
}

% ==================== Code Block Styling ====================
\lstdefinestyle{pythonstyle}{
  language=Python,
  basicstyle=\ttfamily\footnotesize,
  keywordstyle=\color{Navy}\bfseries,
  commentstyle=\color{CodeGreen}\itshape,
  stringstyle=\color{CodePurple},
  backgroundcolor=\color{CodeBg},
  frame=single,
  rulecolor=\color{black!15},
  breaklines=true,
  showstringspaces=false,
  numbers=left,
  numberstyle=\tiny\color{Muted},
  xleftmargin=1.5em,
  framexleftmargin=1.2em
}
\lstset{style=pythonstyle}

% ==================== Header & Footer Setup ====================
\pagestyle{fancy}
\fancyhf{}
\fancyhead[L]{\small\color{Muted}\textsc{TOPIC TITLE}}   % Replace with your topic title
\fancyhead[R]{\small\color{Muted}Revision Notes}
\fancyfoot[C]{\small\color{Muted}\thepage}
\renewcommand{\headrulewidth}{.4pt}

% ==================== Section Headings Styling ====================
\titleformat{\chapter}[display]
  {\bfseries\color{Navy}}{\Large\chaptername\ \thechapter}{4pt}{\Huge}
\titleformat{\section}
  {\Large\bfseries\color{Navy}}{\thesection}{.7em}{}
\titleformat{\subsection}
  {\large\bfseries\color{Blue}}{\thesubsection}{.7em}{}
\titleformat{\subsubsection}
  {\normalsize\bfseries\color{Teal}}{\thesubsubsection}{.7em}{}

% ==================== Helpful Mathematical Notation ====================
\newcommand{\RR}{\mathbb{R}}
\newcommand{\E}{\operatorname{E}}
\newcommand{\Var}{\operatorname{Var}}
\newcommand{\Cov}{\operatorname{Cov}}
\newcommand{\RSS}{\operatorname{RSS}}
\newcommand{\argmin}{\operatorname*{arg\,min}}
\newcommand{\argmax}{\operatorname*{arg\,max}}
\newcommand{\sign}{\operatorname{sign}}
\newcommand{\norm}[1]{\left\lVert #1\right\rVert}
\newcommand{\abs}[1]{\left|#1\right|}

% Vector & Matrix Bold Notation
\newcommand{\X}{\mathbf{X}}
\newcommand{\y}{\mathbf{y}}
\newcommand{\I}{\mathbf{I}}
\newcommand{\bta}{\boldsymbol{\beta}}
\newcommand{\bhat}{\widehat{\boldsymbol{\beta}}}
\newcommand{\eps}{\boldsymbol{\varepsilon}}

% ==================== Document Title & Metadata ====================
\title{\Huge\bfseries [Topic Name: e.g., Logistic Regression]\\[0.2em]}
\author{Noted by Md. Inzamamul Lohani\\
Software Engineering, IIT, DU}
\date{\today}

% ==============================================================================
% DOCUMENT BODY
% ==============================================================================
\begin{document}

% Title Page (matches Ridge & Lasso clean book format)
\hypersetup{pageanchor=false}
\maketitle
\hypersetup{pageanchor=true}

% Frontmatter: Roman page numbering for Table of Contents
\frontmatter
\tableofcontents

% Mainmatter: Arabic page numbering starting at Page 1
\mainmatter

% ==============================================================================
\chapter{The Big Picture \& Motivation}
% ==============================================================================

\section{Introduction and Problem Formulation}
Begin with a clear, approachable overview of the topic. Connect it to concepts the reader already understands.

In multiple linear modeling, we describe the target variable $y \in \RR^n$ via $p$ features $X \in \RR^{n \times p}$:
\begin{equation}
  y = X\beta + \varepsilon, \quad \varepsilon \sim \mathcal{N}(0, \sigma^2 I).
\end{equation}

\begin{keyidea}[title=Core concept]
State the single most fundamental takeaway of this topic right here. Why does this technique exist, and what core problem does it solve?
\end{keyidea}

\section{Geometric and Intuitive Foundation}
Provide the intuition before jumping deep into algebraic manipulations.

\begin{intuition}[title=Visualizing the behavior]
Use intuitive analogies or geometric interpretations. For example, explain why level curves touch constraints at corners or how gradient vectors interact with decision boundaries.
\end{intuition}

\subsection{Comparison Table}
Summarize how this method compares with closely related methods:

\begin{center}
\begin{tabularx}{0.95\textwidth}{>{\bfseries}lXX}
\toprule
Feature / Property & Method A (e.g., Ridge) & Method B (e.g., Lasso) \\
\midrule
Penalty Type & $L_2$ squared norm & $L_1$ absolute norm \\
Solution Type & Analytical closed-form & Iterative (e.g., coordinate descent) \\
Sparsity & Shrinks smoothly towards zero & Exact zeros (feature selection) \\
Behavior with Collinearity & Groups correlated predictors & Picks one arbitrarily \\
\bottomrule
\end{tabularx}
\end{center}


% ==============================================================================
\chapter{Mathematical Foundations \& Derivation}
% ==============================================================================

\section{Objective Function Formulation}
Write the complete optimization problem clearly:
\begin{equation}
  \boxed{J(\beta) = \norm{y - X\beta}_2^2 + \lambda \mathcal{R}(\beta).}
  \label{eq:objective}
\end{equation}
Explain each term carefully:
\begin{itemize}
  \item $\norm{y - X\beta}_2^2$: The data-fidelity (Residual Sum of Squares) term.
  \item $\mathcal{R}(\beta)$: The regularization penalty preventing overfitting.
  \item $\lambda \geq 0$: The hyperparameter balancing fit vs complexity.
\end{itemize}

\section{Step-by-Step Derivation}
Walk through the mathematical steps systematically:

\begin{enumerate}
  \item \textbf{Expand the quadratic loss:}
  \begin{equation}
    \RSS(\beta) = (y - X\beta)^T(y - X\beta) = y^Ty - 2y^TX\beta + \beta^TX^TX\beta.
  \end{equation}

  \item \textbf{Take the matrix gradient with respect to $\beta$:}
  \begin{equation}
    \nabla_\beta \RSS(\beta) = -2X^Ty + 2X^TX\beta.
  \end{equation}

  \item \textbf{Set the gradient to zero and solve:}
  \begin{equation}
    2X^TX\beta - 2X^Ty = 0 \implies \boxed{\bhat = (X^TX)^{-1}X^Ty.}
  \end{equation}
\end{enumerate}

\begin{warningbox}[title=Watch out for singularities]
Remember to mention edge cases, numerical instabilities, or rank deficiency (e.g., when $p > n$ or variables are collinear).
\end{warningbox}


% ==============================================================================
\chapter{Algorithm \& Practical Implementation}
% ==============================================================================

\section{Algorithmic Optimization}
Detail how the numerical optimization is implemented in practice.

\begin{worked}[title=Step-by-step update rule]
Demonstrate a concrete 1D or 2D derivation of the coordinate or gradient update:
\begin{equation}
  \beta_j^{(t+1)} \leftarrow \frac{\operatorname{soft}(x_j^Tr_j,\, \lambda/2)}{x_j^Tx_j}.
\end{equation}
\end{worked}

\section{Python Implementation}
Provide a clean, idiomatic Python code implementation:

\begin{lstlisting}
import numpy as np

def soft_threshold(z: float, threshold: float) -> float:
    """Applies soft thresholding S(z, threshold)."""
    return np.sign(z) * max(abs(z) - threshold, 0.0)

def coordinate_descent_step(X: np.ndarray, y: np.ndarray, beta: np.ndarray, lam: float):
    """Performs one cycle of coordinate descent."""
    n, p = X.shape
    for j in range(p):
        # Compute partial residual excluding feature j
        r_j = y - (X @ beta - X[:, j] * beta[j])
        rho_j = X[:, j] @ r_j
        denom = np.sum(X[:, j]**2)
        beta[j] = soft_threshold(rho_j, lam / 2.0) / denom
    return beta
\end{lstlisting}


% ==============================================================================
\chapter{Practical Considerations \& Tuning}
% ==============================================================================

\section{Feature Standardization}
Explain why feature scaling matters and provide the rule of thumb:
\begin{equation}
  x_{ij}^{(\text{scaled})} = \frac{x_{ij} - \bar{x}_j}{s_j}.
\end{equation}

\section{Hyperparameter Selection via Cross-Validation}
Explain the $K$-fold cross-validation scheme to find $\lambda^*$:
\begin{itemize}
  \item Split training data into $K$ disjoint folds.
  \item Fit model on $K-1$ folds across candidate hyperparameter grid.
  \item Select $\lambda$ minimizing mean validation loss or apply the 1-SE rule.
\end{itemize}


% ==============================================================================
\chapter{Revision Sheet \& Self-Test}
% ==============================================================================

\section{The Essential Formulas}
\begin{align*}
  \text{Cost Function:}\quad& J(\beta) = \RSS(\beta) + \lambda \mathcal{R}(\beta) \\
  \text{Matrix Form:}\quad& \RSS(\beta) = y^Ty - 2y^TX\beta + \beta^TX^TX\beta \\
  \text{Normal Equation:}\quad& X^TX\beta = X^Ty \\
  \text{Optimal Estimator:}\quad& \bhat = (X^TX + \lambda I)^{-1} X^Ty
\end{align*}

\section{Conceptual Summary}
\begin{itemize}
  \item Bullet 1: Core motivation and objective.
  \item Bullet 2: Mechanism of regularization / shrinkage.
  \item Bullet 3: Key trade-offs (e.g., Bias vs. Variance).
  \item Bullet 4: Key operational assumptions and requirements (scaling, independence).
\end{itemize}

\section{Self-Test Questions}
\begin{enumerate}
  \item Why does ordinary least squares fail when predictors are collinear?
  \item What is the geometric difference between $L_1$ and $L_2$ constraint regions?
  \item Why must data be centered/scaled before applying regularization?
  \item When is closed-form calculation preferred over iterative gradient descent?
\end{enumerate}

\section{Short Answers}
\begin{enumerate}
  \item $X^TX$ becomes near-singular, inflating $\Var(\bhat)$.
  \item $L_1$ has sharp corners on axes; $L_2$ is a smooth hypersphere.
  \item Penalties treat all $\beta_j$ equally; different scales would unfairly penalize smaller-unit features.
  \item Closed form is fast for modest $p$; iterative methods scale better to massive datasets ($n, p \gg 10^4$).
\end{enumerate}

\begin{keyidea}[title=Final Mental Model]
Summarize the grand intuition here in 2--3 sentences to cement the concept before moving to the next topic.
\end{keyidea}

\end{document}
